\chapter*{Temple of the Ancients}\stepcounter{chapter}\phantomsection\addcontentsline{toc}{chapter}{Temple of the Ancients}
\DndDropCapLine{H}{\entryfont igh within the Cairngorm Mountains lies a structure whose name has survived far better than its history. The Temple of the Ancients stands isolated among the snow-covered peaks, its true purpose long since forgotten. Only a small portion of the structure remains exposed to the outside world; beyond its monumental entrance, the Temple disappears deep beneath the mountain.}

\section*{First Sight}\stepcounter{section}\phantomsection\addcontentsline{toc}{section}{First Sight}
{\entryfont The final approach leads the party through increasingly barren terrain. Trees become sparse and stunted before disappearing entirely, leaving only dark stone, windblown snow, and the towering peaks of the Cairngorms. Eventually, the party crests a ridge overlooking a shallow mountain basin.}

\begin{DndReadAloud}
	The narrow mountain path climbs one final ridge, and for the first time in hours the land opens before you.

	Across a basin of snow and exposed black stone stands something that has no place among the natural peaks surrounding it. Great blocks of dark stone rise from the mountainside, arranged in severe, geometric shapes untouched by snow or erosion. Above them towers an immense circular arch of stone, its centre filled by a pale golden light that burns steadily despite the howling wind.

	A broad stairway climbs towards the structure and ends before a colossal sealed doorway beneath the glowing ring.

	Yet you are not the first to reach this place.

	At the foot of the stairs, dark shapes interrupt the otherwise untouched snow. Tents. Crates. The remains of a camp.

	And figures are moving among them.
\end{DndReadAloud}

{\noindent\entryfont The Temple itself shows no obvious signs of recent habitation. The camp, however, is recent. From this distance, the party cannot immediately determine who established it or what happened there.

Characters who stop to observe the site before approaching may notice that the figures appear far more interested in the camp and the Temple than in watching the surrounding mountains. This gives cautious characters an opportunity to approach stealthily or investigate from afar before committing themselves.}

\section*{The Ruined Camp}\stepcounter{section}\phantomsection\addcontentsline{toc}{section}{The Ruined Camp}
{\entryfont The camp originally belonged to the Crimson Paw, whose hunters had reached the Temple shortly before servants of Zargothrax discovered them.

Several weather-beaten tents surround a central firepit, alongside stacked travelling crates, hunting equipment, coils of rope, heavy traps, and the remains of provisions. Much of it has since been overturned, searched, or destroyed.

A reinforced iron transport cage stands near the edge of the camp. Inside lies a dead displacer beast, its dark hide pierced by several wounds. Heavy restraints remain fastened around its legs and neck. The creature had evidently been captured alive and was awaiting transport when the camp was attacked.

Among the scattered possessions is a small Rabbit's Foot. Characters familiar with previous Crimson Paw activity may immediately recognise its significance. Otherwise, its presence can serve as another recurring clue connecting this expedition to the poachers.

What the party encounters among the remains depends upon their level.}

\subsection*{Party Below 5th Level}
{\entryfont The Crimson Paw hunters are already dead or missing when the party arrives. Their camp has since been thoroughly ransacked by a small goblinoid force serving Zargothrax.

Use the following creatures:
\begin{itemize}
	\item 1 Hobgoblin Captain
	\item 2 Hobgoblin Iron Shadows
	\item 1 Goblin Boss
	\item 2 Goblins
\end{itemize}
The goblinoids did not establish the camp and are primarily occupied with looting it. Containers have been broken open, equipment scattered, and anything apparently valuable gathered for later removal.

As most of the group is distracted by the camp, a stealthy approach is a particularly viable option and may allow the party to scout the encounter or gain an advantageous position before being noticed.}

\subsection*{Party of 5th Level or Higher}
{\entryfont The creatures occupying the camp are not looters. They are part of a deliberate expedition sent by Zargothrax.

Use the following creatures:
\begin{itemize}
	\item 1 Bone Naga
	\item 1 Cultist of Zargothrax
	\item 2 Sword Wraith Warriors
\end{itemize}
The Cultist of Zargothrax is actively investigating the Temple entrance, using ritual components and arcane implements in an attempt to determine how it can be opened.

Meanwhile, the Bone Naga and Sword Wraith Warriors stand guard, watching the camp and surrounding approaches for intruders. Consequently, approaching the camp unnoticed is considerably more difficult than in the lower-level encounter.

None of the cultist's attempts to open the Temple have succeeded.}

\begingroup
	\DndSetThemeColor[PhbTan]
	\begin{DndComment}{DM Annotation - Heart Pure of Steel}
		The Heart Pure of Steel is required to open the entrance to the Temple of the Ancients. If the party has not previously obtained it as a reward from Lady Iona McDougall or Mr. Owl Hooters, they find it among the remains of Zargothrax underlings.

		This serves as a fail-safe to ensure that the party can always progress into the Temple. If they already possess the Heart Pure of Steel, no additional copy is found here.
	\end{DndComment}
\endgroup

\section*{The Ancient Gate}\stepcounter{section}\phantomsection\addcontentsline{toc}{section}{The Ancient Gate}
{\entryfont A broad flight of stone steps leads from the camp towards the Temple's entrance. The enormous doorway has no visible handle, hinges, or conventional lock. At its centre lies a small recess whose shape closely resembles the shape of a heart.

As a character carrying the Heart Pure of Steel approaches the entrance, the artifact gradually becomes warm. The sensation intensifies with every step towards the door, making the connection between the Heart and the Temple increasingly apparent.

Despite the matching recess, the Heart does not react further unless the Autumnal Equinox has arrived.}

\begingroup
	\DndSetThemeColor[PhbTan]
	\begin{DndComment}{DM Annotation - The Autumnal Equinox}
		The Temple can only be opened during the Autumnal Equinox, occurring between \textbf{22nd and 24th September 992 AD}. The campaign begins on 13th September, leaving the exact date within this range for the DM to choose according to the party's progress.

If the party reaches the Temple before this date, the increasing warmth of the Heart Pure of Steel confirms that they have found the correct location, but the entrance remains sealed. The party must wait until the chosen date before proceeding.
	\end{DndComment}
\endgroup

{\noindent\entryfont When the party approaches the entrance on the chosen day, read or paraphrase:}

\begin{DndReadAloud}
	With each step towards the towering door, the Heart grows warmer—until suddenly it tears itself from your grasp.

	Suspended in the air, it drifts towards the stone and settles perfectly into the small recess at its centre.

	For a moment, nothing happens.

	Then a deep vibration passes through the mountain.

	Ancient lines carved into the doorway flare with pale golden light. Stone grinds against stone as mechanisms dormant for centuries awaken somewhere beyond the wall. Slowly, the colossal door retreats into the mountain, revealing a dark passage beyond.

	The Heart pulls free from the stone and flies back into your waiting hand.

	The way into the Temple of the Ancients stands open.
\end{DndReadAloud}

{\noindent\entryfont Once opened, the entrance remains accessible and the Heart Pure of Steel returns to the character who previously carried it. It is not consumed by the process.}

\section*{Dungeon}\stepcounter{section}\phantomsection\addcontentsline{toc}{section}{Dungeon}\subsection*{\inTextCircled[10]{A}{10} Entrance}
{\entryfont Beyond the entrance, a broad stone staircase descends deep into the mountain. The steps continue far enough that the daylight from the open gate gradually dwindles behind the party, leaving the passage increasingly dependent upon whatever light they carry with them.

The air within is cold and remarkably still. Dust lies thick upon the stairs, undisturbed by footprints or other signs of recent passage. Despite the Temple's immense age, the stonework remains unusually well preserved.}

\subsection*{Loot Tables}
\subsubsection*{Minor Treasures}
\begin{DndTable}[header=Random Minor Treasure]{lX}
	d12	& Discovery\\
	1	& 2d6 x 10 GP in ancient coins\\
	2	& 3d6 x 10 SP and 1d6 x 10 GP\\
	3	& A small gemstone worth 50 GP\\
	4	& 1d4 gemstones worth 25 GP each\\
	5	& An ornate ceremonial vessel worth 75 GP\\
	6	& A set of ancient silver implements worth 100 GP\\
	7	& 1d4 pieces of finely preserved incense worth 25 GP each\\
	8	& A sealed packet containing Spell Scroll (1st Level)\\
	9	& A sealed packet containing Spell Scroll (2nd Level)\\
	10	& One roll on the Random Potion Table\\
	11	& An ancient object of unclear purpose worth 100 GP to a collector\\
	12	& Nothing\\
\end{DndTable}
\subsubsection*{Potions}
\begin{DndTable}[header=Random Potion Table]{lXX}
	d20		& Potion						& Marking									\\
	1 - 3	& Potion of Healing				& Clearly marked							\\
	4 - 5	& Potion of Climbing			& Clearly marked							\\
	6 - 7	& Potion of Animal Friendship	& Faded marking								\\
	8 - 9	& Potion of Water Breathing		& Clearly marked							\\
	10 - 11	& Potion of Growth				& Unmarked									\\
	12 - 13	& Potion of Resistance			& Faded marking; roll damage type randomly	\\
	14		& Potion of Greater Healing		& Clearly marked							\\
	15		& Potion of Fire Breath			& Unmarked									\\
	16		& Potion of Heroism				& Faded marking								\\
	17		& Potion of Diminution			& Unmarked									\\
	18		& Potion of Mind Reading		& Unmarked									\\
	19		& Potion of Invisibility		& Unmarked									\\
	20		& Roll twice, ignoring further results of 20	& —
\end{DndTable}

\subsection*{\inTextCircled[10]{1}{10} The Lower Hall}
{The narrow entrance opens into a broad, angular hall stretching through the lower reaches of the Temple. Raised sections, alcoves and adjoining chambers interrupt its otherwise imposing symmetry, while another staircase at its far end descends deeper beneath the mountain.}
\subsubsection*{Hidden Passage}
{\entryfont Along the northern wall, one section of masonry conceals an exceptionally narrow opening. A character examining the area can discover the passage with a successful \textbf{DC 15 Wisdom (Perception) Check}.

The opening is far too narrow for most creatures to traverse. Only a Small or smaller creature can squeeze through it. A creature capable of moving through extremely confined spaces, such as one under the effects of Gaseous Form, can also pass through.}

\subsection*{\inTextCircled[10]{2}{10} Minor Treasure Room}
{\entryfont This small chamber appears to have once served as a repository for valuables and ceremonial objects. Low stone shelves and several partially decayed containers line the walls, their contents covered beneath centuries of dust.

Most of what remains has little value beyond its age, but a careful search reveals a handful of objects that have survived in good condition.

Roll 1d2 times on the Random Minor Loot Table to determine what the party finds. Any coins, jewellery, or crafted objects discovered here are distinctly ancient in design, though their exact origin is difficult to determine.}

\subsection*{\inTextCircled[10]{3}{10} Potion Storage}
{\entryfont Stone shelves line the walls of this narrow chamber, holding dozens of bottles, flasks, and stoppered vessels. Many have cracked or dried out over the centuries, while others remain remarkably well preserved. Some still bear faded labels or simple alchemical symbols, whereas others are entirely unmarked.

Searching the shelves reveals a Potion of Gaseous Form alongside 1d4 additional intact potions. Roll once on the Random Potion Table for each additional potion found.}

\subsection*{\inTextCircled[10]{4}{10} Hidden Passage}
{\entryfont The narrow opening leads into a cramped passage that extends surprisingly far into the Temple. Although its entrance is barely large enough for a Small creature to pass, portions of the passage beyond widen considerably before narrowing again.

Several skeletons lie scattered along its length. Most belonged to Medium or larger creatures and bear no obvious signs of violence. Some are wedged awkwardly into particularly narrow sections, while others lie beside the walls where they apparently died after becoming trapped within the passage. How they originally entered remains unclear.

Characters searching the remains find little of value, though the unusual collection of skeletons makes it apparent that this hidden route has claimed more than one unfortunate explorer.

At the far end of the passage sits another skeleton unlike the others. Its remains are still clad in a suit of remarkably clean, shining armour, seemingly untouched by the centuries that have reduced everything around it to dust and bone.

If a character touches either the skeleton or its armour, they immediately disappear from the passage and are teleported elsewhere within the Temple.}

\begingroup
	\DndSetThemeColor[PhbTan]
	\begin{DndComment}{DM Annotation - Teleportation}
		The teleportation occurs immediately upon physical contact and requires no saving throw. Only the creature touching the skeleton is transported; other creatures within the passage remain behind.

Continue with the corresponding room when this occurs.
	\end{DndComment}
\endgroup

\subsection*{\inTextCircled[10]{5}{10} Abandoned Stores}
{\entryfont The hallway continues until it terminates at a collection of old crates, boxes and reinforced chests stacked against the walls. Most have deteriorated considerably, though several stone and metal containers remain intact.

The majority contain mundane supplies rendered useless by age: fragments of cloth, corroded tools, empty vessels and unidentifiable remnants of other materials once stored here.

A thorough search of the containers reveals one roll on the Random Minor Loot Table. In addition, the party finds a set of three intact flasks of Holy Water, carefully sealed inside a padded metal case.

One reinforced chest remains locked. It can be opened with a successful DC 14 Dexterity check using Thieves' Tools or forced open with a successful DC 16 Strength (Athletics) check. Inside are 2d6 x 10 GP in ancient coinage and a small gemstone worth 50 GP.}

\subsection*{\inTextCircled[10]{6}{10} Scrying Chamber}
{\entryfont A heavy stone door seals this side chamber. Unlike most doors within the Temple, it has no obvious mechanism for opening it from the hallway. The door can be forced open with a successful DC 20 Strength (Athletics) check. Alternatively, the DM may allow sufficiently destructive magic or appropriate tools to breach it.

Beyond lies a small chamber filled with ancient scrying equipment. Tarnished metal frames surround cloudy crystals, shallow stone basins stand upon pedestals, and faded arcane diagrams cover portions of the walls. None of the equipment appears immediately functional.

Dominating the chamber is a large map depicting much of the Kingdom of Fife and the surrounding lands. Several locations have been deliberately marked:
\begin{itemize}
	\item Swamp of Paisley
	\item Loch Leven
	\item Aberdeen
	\item Island of Unst
	\item Auchtertool
	\item Temple of the Ancients
	\item an otherwise unmarked location within the Cairngorm Mountains, lying between the Temple and Aberdeen
\end{itemize}}

\subsection*{\inTextCircled[10]{7}{10} Boobytrapped Hallway}
{\entryfont This long, winding hallway is protected by a series of mechanical and magical traps. Unlike the deterioration found elsewhere in the Temple, the mechanisms here remain functional despite their immense age.

The traps are spread throughout the hallway rather than concentrated in a single location. A character actively searching for traps can detect a trap with a successful \textbf{DC 14 Wisdom (Perception) Check}. Once detected, the party can attempt to disable or otherwise circumvent it.}
\begin{DndTrap}[width=0.5\textwidth - 4pt]{Pressure-Plate Darts}
	\DndTrapBasics[
		save_dc			= 13,
		skill			= {Dexterity Saving Throw},
		pos_reaction	= {Duck},
		neg_reaction	= {Brace},
		damage			= {1d6 Piercing + 1d4 Poison},
	]
	{\entryfont\noindent A concealed pressure plate causes small openings along both walls to release a volley of ancient metal darts. On a failed saving throw, the creature takes the listed damage, or half as much on a successful one.}
\end{DndTrap}
\begin{DndTrap}[width=0.5\textwidth - 4pt]{Runic Flame Jet}
	\DndTrapBasics[
		save_dc			= 14,
		skill			= {Dexterity Saving Throw},
		pos_reaction	= {Step Back},
		neg_reaction	= {Duck},
		damage			= {2d6 Fire},
	]
	{\entryfont\noindent Stepping across a nearly invisible rune causes ancient channels within the walls to ignite, releasing a sudden blast of flame across the passage. A creature failing the saving throw takes the listed damage, or half as much on a successful one.}
\end{DndTrap}
\begin{DndTrap}[width=0.5\textwidth - 4pt]{Pendulum Blade}
	\DndTrapBasics[
		save_dc			= 15,
		skill			= {Dexterity Saving Throw},
		pos_reaction	= {Duck},
		neg_reaction	= {Jump},
		damage			= {2d8 Slashing},
	]
	{\entryfont\noindent A concealed mechanism releases a heavy curved blade that sweeps across the hallway at roughly chest height. A creature failing the saving throw takes the listed damage, or half as much on a successful one.}
\end{DndTrap}

\subsection*{\inTextCircled[10]{8}{10} Forgotten Storage}
{\entryfont A small storage chamber branches from the trapped hallway. Several old containers remain stacked against the walls, most of their contents having deteriorated beyond recognition.

Searching the surviving containers grants one roll on the Random Minor Loot Table.

A successful DC 17 Wisdom (Perception) check reveals an unusually narrow seam behind one of the shelving units. Moving the shelf aside exposes a concealed passage leading directly into the adjacent chamber. The passage appears to have been intentionally hidden and can be traversed normally once discovered.}

\subsection*{\inTextCircled[10]{9}{10} Sleeping Quarters}
{\entryfont This secluded chamber contains three stone-framed beds arranged along its walls. Despite the age of everything surrounding them, each bed still bears a thin mattress and faded coverings that remain strangely intact.

Nothing distinguishes one bed from another at first glance.

Whenever a creature lies down upon one of the beds, their surroundings immediately fade from perception. Their body remains motionless upon the bed while their consciousness experiences a vivid vision determined by the bed they chose.

The three beds produce different visions:}
\subsubsection*{\inTextCircled[10]{9A}{10} Bed of Nightmares}
{\entryfont The character experiences a terrifying vision shaped around their deepest fears, destruction, and impending loss. At the end of the vision, they must succeed on a DC 14 Wisdom saving throw, taking 3d6 Psychic damage on a failed save, or half as much damage on a successful one.}
\begin{DndReadAloud}
	You awaken violently, gasping for breath as the lingering pain of the vision remains painfully real.
\end{DndReadAloud}
\subsubsection*{\inTextCircled[10]{9B}{10} Bed of Destinies (True)}
{\entryfont The character witnesses a fragment of events that may genuinely come to pass.}
\subsubsection*{\inTextCircled[10]{9C}{10} Bed of Destinies (False)}
{\entryfont The character receives an equally convincing vision of a future that will never occur.}

{\entryfont The visions last only moments in the outside world. When they end, the character awakens immediately and suffers no inherent physical harm.}

\subsection*{\inTextCircled[10]{9}{10} Ancient Smithy}
{The rhythmic ringing of metal can be heard long before this chamber comes into view. Within, an ancient forge continues to operate without anyone tending it. Hammers, tongs, chisels, and other smithing tools fly between anvils and workbenches as though wielded by invisible hands.

The animated tools ignore creatures entering the chamber, though characters should be wary of their erratic movements. A carelessly placed head may quickly find itself in the path of a passing hammer.

At the centre of the Smithy stands a large stone anvil covered in faintly glowing runes. Placing a piece of protective equipment - such as armor, a shield, coat, cloak, or similar garment - upon the anvil causes the forge to immediately spring into action. Weapons and other objects are briefly inspected by the tools before being unceremoniously pushed from the anvil.

When the Smithy springs to life, its tools suddenly tear themselves from their resting places and hurtle towards the anvil. Each creature within the chamber must make a DC 13 Dexterity saving throw as hammers, tongs, and other implements fly past them. On a failed save, a creature takes 1d6 Bludgeoning damage from a wayward tool. On a successful save, the creature avoids the flying implements unharmed.

The Smithy can enhance both mundane and magical equipment. Once the process begins, roll on the appropriate Smithy Outcome Table. If successful, the item retains all of its existing properties and gains the rolled enchantment.}

\begin{DndTable}[header=Mundane Equipment]{lX}
	d20		& Outcome\\
	1 - 2	& The item is destroyed.\\
	3 - 17	& Roll on the Uncommon Smithy Enchantments table.\\
	18 - 20	& Roll on the Rare Smithy Enchantments table.\\
\end{DndTable}
\begin{DndTable}[header=Magical Equipment]{lX}
	d20		& Outcome\\
	1 - 8	& The item and its existing magic are destroyed.\\
	9 - 17	& Roll on the Uncommon Smithy Enchantments table.\\
	18 - 20	& Roll on the Rare Smithy Enchantments table.\\
\end{DndTable}

{\noindent\entryfont The destruction of an item is permanent. Once the tools begin working, the item cannot safely be removed from the anvil until the process has concluded.}

\begin{DndTable}[header=Uncommon Smithy Enchantment]{llX}
	d8	& Enchantment		& Effect\\
	1	& Fleet				& While wearing or carrying the item, the bearer's Speed increases by 5 feet.\\
	2	& Surefooted		& The bearer has Advantage on saving throws against being knocked Prone or moved against their will.\\
	3	& Watchful			& The bearer cannot be Surprised while conscious.\\
	4	& Steadfast			& The bearer has Advantage on saving throws against being Frightened.\\
	5	& Featherbound		& The bearer can cast Feather Fall on themselves once per Long Rest without requiring material components.\\
	6	& Guardian			& While wearing or carrying the item, the bearer gains a +2 bonus to Initiative.\\
	7	& Unyielding		& Once per Long Rest when the bearer would be reduced to 0 Hit Points but not killed outright, they are reduced to 1 Hit Point instead.\\
	8	& Element-Touched	& Randomly determine Acid, Cold, Fire, Lightning, or Thunder. Once per Long Rest, the bearer can gain Resistance to that damage type for 1 minute as a Bonus Action.\\
\end{DndTable}
\begin{DndTable}[header=Rare Smithy Enchantment]{llX}
	d8	& Enchantment		& Effect\\
	1	& Elemental Ward		& Randomly determine Acid, Cold, Fire, Lightning, or Thunder. The bearer gains Resistance to that damage type.\\
	2	& Spellguard		& The bearer has Advantage on saving throws against spells and other magical effects.\\
	3	& Resolute			& Once per Long Rest when the bearer fails a saving throw, they can reroll it, taking the new result.\\
	4	& Blinking			& As a Bonus Action, the bearer can teleport up to 15 feet to an unoccupied space they can see. This property can be used once per Short or Long Rest.\\
	5	& Ghostforged		& Once per Long Rest, the bearer can become partially incorporeal until the end of their turn, allowing them to move through creatures and objects as Difficult Terrain. If they end their turn inside an object, they take 1d10 Force damage and are shunted to the nearest unoccupied space.\\
	6	& Defiant			& When the bearer takes damage, they can use their Reaction to gain Resistance to that instance of damage. This property can be used once per Long Rest.\\
	7	& Arcane Bulwark	& Once per Long Rest when the bearer is targeted by a spell, they can use their Reaction to gain Advantage on any saving throw caused by that spell and Resistance to its damage until the end of the current turn.\\
	8	& Second Skin		& As a Bonus Action, the bearer can magically don or doff the item. While worn, the item can also alter its outward appearance according to the bearer's wishes without changing its statistics.
\end{DndTable}

\begingroup
	\DndSetThemeColor[PhbTan]
	\begin{DndComment}{DM Annotation - Smithy's Price}
		Each character can offer only one item to the Ancient Smithy. After it has accepted an item from a character, the tools refuse anything else subsequently offered by that character.

	The increased risk when enhancing an existing magic item is intentional. The Smithy is capable of weaving additional magic into an already enchanted object, but doing so is considerably less stable.

	Artifacts, Vestiges of Inverness, and other plot-critical items are never accepted by the Smithy. The tools simply return such an object without attempting to alter it.
	\end{DndComment}
\endgroup

\subsection*{\inTextCircled[10]{10}{10} Secluded Workshop}
{\entryfont Several shelves filled with deteriorated tools and unidentifiable materials line this otherwise isolated chamber.

Against one wall rests a skeleton clad in the same immaculate, shining armour as the skeleton at the opposite end of the hidden passage. Touching either the skeleton or its armour immediately teleports the character back to its counterpart.

No saving throw is required, and only the creature making physical contact is transported.}

\subsection*{\inTextCircled[10]{11}{10} Ancient Kitchen}
{\entryfont Long stone counters, ovens, cupboards, and preparation tables fill this remarkably well-preserved kitchen. Despite having remained sealed for centuries, platters of food and sealed vessels of drink remain fresh. The meals are warm, the fruit unspoiled, and the drinks pleasantly cool, as though everything had been prepared only moments before the party arrived.

The food is exceptionally nourishing, and the chamber provides a safe place to recover.

A Short Rest completed within the Kitchen grants the benefits of two Short Rests. A character can therefore spend Hit Dice and use features that recover during a Short Rest twice.

In addition, a character with the Spellcasting or Pact Magic feature can recover expended spell slots with a combined level equal to half their Proficiency Bonus, rounded up. None of the recovered spell slots can be higher than 5th level.

After a character has benefited from the Kitchen in this way, they cannot gain its special benefits again until they finish a Long Rest.}

\subsection*{\inTextCircled[10]{12}{10} Hall of Stories}
{\entryfont The walls of this broad chamber are covered almost entirely by enormous tapestries. Though their colours have faded with age, the fabric itself remains remarkably intact. Some depict recognisable events from the earliest histories of Dundee and the Kingdom of Fife, while others portray scenes for which no surviving account appears to exist.

Characters who take the time to examine the tapestries can discover the following scenes.}

\subsubsection*{The Founding of Dundee}
{\entryfont Several interconnected tapestries depict the earliest days of Dundee. They show the first settlement growing beneath the protection of its founders, its walls rising, and scattered peoples gathering beneath a common banner. Later scenes show Dundee expanding beyond its beginnings and becoming the heart around which the Kingdom of Fife would eventually form.

Although many of the individuals depicted are recognisable from surviving legends, several scenes differ subtly from their modern retellings. Whether these differences represent forgotten history or merely the interpretation of whoever created the tapestries is impossible to determine.}

\subsubsection*{The Three Relics}
{\entryfont One tapestry is dominated not by people or places, but by three objects presented with unusual prominence.

The first is a mighty warhammer of elaborate construction, immediately recognisable from countless legends as the Hammer of Glory.

Beside it stands a long, wickedly shaped blade covered in cruel ornamentation: the legendary Knife of Evil.

The third object is considerably harder to understand. It resembles neither weapon nor tool known within Fife: an angular construction of metal plates surrounding tubes, cylinders, small levers, glass-like components, and other mechanisms of completely unknown purpose. One end appears designed to be held, while the opposite terminates in several hollow metal openings. No bowstring, blade, striking surface, or other obvious means of using it as a weapon can be seen.

Together, the three objects are presented with a reverence afforded to almost nothing else within the Hall.}

\subsubsection*{The Dragon and the Talisman}
{\entryfont Another tapestry depicts an immense dragon towering over a devastated landscape, its wings spread across much of the woven scene.

Opposite the creature hangs a small circular talisman, depicted disproportionately large so that its intricate construction can be clearly seen. A brilliant light radiates from its centre toward the dragon.

No accompanying scene reveals the outcome of their confrontation, nor does the tapestry explain who carried the talisman or why it was used against the creature.}

\subsubsection*{The Trial of Elements}
{\entryfont A group of warriors battles strange beings formed from flame, stone, storm, and other elemental forces. Swords, spears, and arrows strike the creatures, yet the force of each blow appears to flow harmlessly into their bodies. In subsequent images, that same energy is unleashed back upon the attacking warriors.

Elsewhere within the tapestry, those same creatures are struck by fire, frost, lightning, acid, and other supernatural forces. Here, the beings recoil and break apart beneath the assault.}

\subsubsection*{The Trial of Might}
{\entryfont Heavily armed warriors facing a brutal host of monstrous combatants.

Flames engulf them. Lightning courses across their bodies. Waves of supernatural energy crash against their ranks. Yet each attack seems to strengthen rather than injure them, the absorbed power gathering within their weapons before being violently returned against their attackers.

In contrast, blades cut deeply, arrows pierce armour, and crushing blows send the creatures tumbling to the ground.}

\subsubsection*{The Trial of Death}
{\entryfont Robed figures surround an enormous undead creature standing amid a field of fallen warriors. The dead appear not merely to lie around the creature, but within it: pale silhouettes of warriors can be seen beneath its translucent form, seemingly trapped inside.

In the following scene, more warriors have fallen and more spectral figures fill the creature. Yet rather than struggling against it, the spirits now appear to move together.

In the final image, the undead monstrosity has turned away from the remaining warriors.

It faces the robed figures that summoned it.}

\begingroup
	\DndSetThemeColor[PhbTan]
	\begin{DndComment}{DM Annotation - Foreshadowing the Three Trials}
		The final three tapestries foreshadow the mechanics of the upcoming Challenge of Elements, Challenge of Might, and Challenge of Death. They should provide clues rather than explicit instructions.

Do not explain their meaning to the players. If they remember the imagery once the trials begin, they should be allowed to draw their own conclusions from it. The effects of attacks during the trials should be described clearly enough that the connection can be recognised.
	\end{DndComment}
\endgroup

\subsection*{\inTextCircled[10]{13}{10} Hall of Wells}
{\entryfont The sound of gently flowing water fills this chamber. Several stone fountains and shallow wells are arranged around its walls, each fed by an unseen source somewhere within the mountain. Despite the Temple's age, their waters remain perfectly clear.

Four of the fountains radiate faint magic. Drinking from one grants the creature its associated blessing:}
\begin{DndTable}[header=Fountain's Blessing]{lX}
	Fountain			& Blessing\\
	Well of Vigor		& Gain 2d6 + 4 Temporary Hit Points until the next Long Rest.\\
	Well of Clarity		& Gain Advantage on the next Intelligence, Wisdom, or Charisma saving throw made before the next Long Rest.\\
	Well of Swiftness	& The creature's Speed increases by 10 feet for 1 hour.\\
	Well of Fortune		& Gain 1d6, which can be added to one ability check or saving throw made before the next Long Rest. The die can be added after the d20 is rolled, but before the outcome is known.\\
\end{DndTable}
{\entryfont A creature can benefit from only one fountain. Once it drinks from one of the four, the water of the remaining fountains tastes like ordinary mountain water and grants that creature no magical effect.

The fountains themselves provide no indication of their exact effects. However, subtle differences in their construction and the symbols carved around each basin may provide clues to the nature of their respective blessings.}

\begingroup
	\DndSetThemeColor[PhbTan]
	\begin{DndComment}{DM Annotation - Choosing a Fountain}
		The fountains are intended to present a minor choice rather than a puzzle with a single correct solution. Allow characters who examine the carvings to make a \textbf{DC 13 Intelligence (Arcana) Check} to gain a general impression of a fountain's effect, such as vitality, clarity, speed, or fortune.

	Drinking from multiple fountains provides no additional benefit, and water removed from the chamber loses its magical properties.
	\end{DndComment}
\endgroup

\subsection*{\inTextCircled[10]{14}{10} Chamber of the Heroes of Old}
{\entryfont Rows of stone sarcophagi fill this solemn chamber, each bearing the carved likeness of a warrior laid to rest within. Names and inscriptions cover their sides, though many have become too weathered to read. These are the resting places of heroes from ages long past, most of whom have disappeared entirely from modern histories.

Resting atop each sarcophagus is a weapon once carried by the hero buried beneath it. Unlike the surrounding stone and faded inscriptions, these weapons remain in perfect condition. Blades retain flawless edges, polished metal reflects the light, and not a trace of rust or decay can be found upon them.

An inscription carved prominently into the chamber reads:}
\begin{DndReadAloud}
	Their deeds are remembered.\\
	Their burdens are ended.\\
	Let the arms of the fallen rest with those who bore them.
\end{DndReadAloud}
{\noindent\entryfont The weapons are not physically secured and can be removed freely. Each functions as a +2 Weapon of its respective type while within the Temple.

Nothing immediately happens when a weapon is taken. No guardian appears, no trap is triggered, and the character carrying it notices no curse or other magical effect.

However, taking the possessions of the interred heroes carries consequences during the trials ahead.}

\begingroup
	\DndSetThemeColor[PhbTan]
	\begin{DndComment}{DM Annotation - Burden of the Dishonoured}
		A character who removes a weapon from a sarcophagus and carries it into the Ritual Chamber gains the Burden of the Dishonoured for each weapon taken.

For each Burden, the character suffers a cumulative -2 penalty to all saving throws made during the Three Trials.

A character can avoid this penalty by returning the weapon to its proper resting place before the ritual begins. Once the ritual has been activated, returning it is no longer possible until the trials are complete.

Do not inform the players of the Burden when a weapon is taken. The inscription and unusual preservation of the weapons serve as their warning.
	\end{DndComment}
\endgroup

{\noindent\entryfont The weapons retain their +2 enchantment throughout the Three Trials and can therefore provide considerable immediate power despite their associated penalty.

Once removed from the Temple of the Ancients, however, any stolen weapon rapidly deteriorates. Its magical properties disappear, polished surfaces tarnish, and centuries of suppressed decay take hold within moments. What remains is an ordinary, badly rusted weapon with no magical properties or significant monetary value.}

\subsection*{\inTextCircled[10]{15}{10} Ritual Chamber}
{\entryfont At the heart of the Temple lies a large circular chamber. Its walls are largely bare, drawing all attention towards a raised stone platform at its centre.

Upon the platform stands a round pedestal engraved with an enormous Crest of Dundee. The same pattern is repeated throughout the surrounding floor and architecture, mirroring the unusual shape of the Temple itself.

At the centre of the crest is a small recess.

As the Heart Pure of Steel is brought into the chamber, its appearance begins to change. Its previous form slowly folds and shifts until the artifact resembles a small talisman bearing the Crest of Dundee, perfectly matching the recess within the pedestal.

Nothing else within the chamber reacts until the Heart is placed inside.

Beginning the Ritual

When the Heart Pure of Steel is placed into the recess, the ritual begins immediately.

The Heart locks firmly into place as pale light spreads outward through the engraved Crest. Lines previously indistinguishable from the surrounding stone illuminate across the pedestal, floor, walls, and ceiling until the entire chamber is traced by the same glow.

The entrance seals behind the party.

The Three Trials then begin in succession:
\begin{itemize}
	\item[1.] Challenge of Elements
	\item[2.] Challenge of Might
	\item[3.] Challenge of Death
\end{itemize}

The party does not physically travel to separate chambers for these trials. Instead, the Ritual Chamber transforms around them. Walls disappear behind supernatural landscapes, the floor changes beneath their feet, and creatures manifest as though suddenly made real.

Once a trial is completed, its surroundings dissolve and the chamber briefly returns before the next begins.}

\begingroup
	\DndSetThemeColor[PhbTan]
	\begin{DndComment}{DM Annotation - The Three Trials}
		The Three Trials are as much puzzles as combat encounters. Each follows a particular rule foreshadowed by the tapestries within the Hall of Stories.

Do not directly reveal these mechanics. Instead, clearly describe how ineffective attacks are absorbed and subsequently used by the creatures. Players who recognise the pattern and adapt their tactics should gain a significant advantage.

Characters carrying weapons stolen from the Chamber of the Heroes of Old suffer the effects of their Burden of the Dishonoured throughout all three trials.
	\end{DndComment}
\endgroup

\subsubsection*{Challenge of Elements}
{\entryfont As the first trial begins, the Ritual Chamber disappears behind a maelstrom of elemental energy. Flames race across broken stone, lightning tears through the air, and clouds of frost and corrosive mist gather around the battlefield.

From this chaos emerge the creatures of the trial.

The exact creatures used for the encounter can be adjusted to suit the party's level. Elementals and creatures capable of dealing elemental damage are particularly appropriate.

All hostile creatures in this trial gain the following properties:

\paragraph*{Physical Absorption} The creature is immune to Bludgeoning, Piercing, and Slashing damage. Whenever it would take damage of one of these types, it instead absorbs the damage, up to a maximum of 20 points at a time.

The next time the creature deals damage with an attack, it expends all stored energy and deals additional Psychic damage equal to the amount stored.

\paragraph*{Elemental Vulnerability} The creature is vulnerable to Acid, Cold, Fire, Lightning, Poison, and Thunder damage.

Whenever physical damage is absorbed, describe the impact visibly flowing into the creature rather than simply telling the player that the attack dealt no damage. The stored energy should become increasingly apparent until it is released.

Once every hostile creature has been defeated, the elemental landscape collapses into pale light and the Ritual Chamber briefly returns.}

\subsubsection*{Challenge of Might}
{\entryfont The chamber transforms again. The elemental chaos disappears, replaced by a stark arena of stone. Heavy footsteps echo across it as brutal warriors materialise around the party.

This trial pits the party against Tanarukks and Gladiators. Adjust their number to provide an appropriate encounter for the party's level.

All hostile creatures in this trial gain the following properties:

\paragraph*{Elemental Absorption} The creature is immune to Acid, Cold, Fire, Lightning, Poison, and Thunder damage. Whenever it would take damage of one of these types, it instead absorbs the damage, up to a maximum of 20 points at a time.

The next time the creature deals damage with an attack, it expends all stored energy and deals additional Force damage equal to the amount stored.

\paragraph*{Physical Vulnerability} The creature is vulnerable to Bludgeoning, Piercing, and Slashing damage.

The visual effect mirrors the previous trial. Elemental attacks disappear into the warriors rather than harming them, leaving supernatural energy coursing across their bodies and weapons until it is discharged with their next attack.

Once all hostile creatures have been defeated, the arena fades and the party returns briefly to the Ritual Chamber.}

\subsubsection*{Challenge of Death}
{\entryfont For the final time, the Ritual Chamber disappears.

The party finds themselves standing upon dead earth beneath a lightless sky. Gravestones and broken monuments surround them while robed cultists complete a ritual around an enormous mass of corpses and spectral figures.

A Graveyard Revenant rises among them.

Use a number of Death Cultists appropriate to the size of the party. Each Cultist retains the abilities of its normal stat block but has only 1 Hit Point.

The following special rules apply throughout this trial:

\paragraph*{Death's Dominion} The Graveyard Revenant and Death Cultists are immune to all damage originating from a creature that is not Undead. This immunity cannot be ignored or overcome by changing damage type.

\paragraph*{Death Within the Trial} A character reduced to 0 Hit Points does not fall Unconscious and makes no Death Saving Throws. Instead, their body dissolves into spectral light and is drawn into the Graveyard Revenant. For the remainder of the trial, that character is considered Undead and exists within the Revenant.

A character cannot truly die while participating in this trial.

\paragraph*{Dominating the Revenant} At the beginning of each of the Graveyard Revenant's turns, every player character currently trapped within it makes a DC 15 Wisdom saving throw.

If at least one character succeeds, the trapped characters seize control of the Revenant for that turn. The players collectively determine its movement, actions, and targets.

If every character fails, the Revenant acts normally under the DM's control.

Because the Death Cultists possess only 1 Hit Point, any successful damaging attack made against one by the Revenant is generally sufficient to kill it.

As more characters are defeated and absorbed, more Wisdom saving throws are made each round, progressively increasing the party's ability to wrest control of the creature.}

\subsubsection*{Completing the Ritual}
{\entryfont Once the Challenge of Death has been completed, every manifested creature and supernatural alteration disappears. The party finds themselves once again standing around the central pedestal.

For several moments, nothing happens.

Then a crack appears across the Heart Pure of Steel.

Read or paraphrase:}
\begin{DndReadAloud}
	A sharp crack breaks the silence.

	A fracture runs through the glowing talisman. Then another. And another.

	The Heart breaks apart - but none of its pieces fall.

	One fragment rises for each of you, suspended in the air above the pedestal. Slowly, the pieces drift outward until a single fragment hangs before each of your faces, glowing with a pale golden light.

	The chamber falls utterly still.

	Then the fragments flare.

	For a single heartbeat, there is nothing but light.

	And the Temple disappears.
\end{DndReadAloud}
{\noindent\entryfont Each character is immediately transported alone to their respective Hall of the Vestige. The fragments of the Heart Pure of Steel accompany them.

The Ritual Chamber itself remains within the Temple, but the Halls to which the characters are transported do not occupy any physical location shown within its ordinary architecture.}

\begingroup
	\DndSetThemeColor[PhbTan]
	\begin{DndComment}{DM Annotation - The Halls of the Vestiges}
		From this point onward, resolve each character's Hall individually. The Halls are bespoke spaces intended to reflect both the character and the Vestige awaiting them. They are described separately in Halls of the Vestiges.

The number of Halls is not determined by the dungeon map. There is always one Hall for each participating player character.
	\end{DndComment}
\endgroup